% =====================================================================
% Note 9, decoding a latent direction. Moved out of the Supplementary Methods, 24 November (guide's
% comment: no Supplementary Methods; NMI allows algorithm detail as a Note).
% =====================================================================

\section{Decoding a latent direction}
\label{note:decoding}

The decoder turns a latent position into expected counts, so moving a cell
along a latent direction and decoding it again shows which genes the
direction carries. \ref{note:claims} uses this to read the immune disease
direction as genes. In this note we give the procedure, the two checks that
had to pass before any ranking, and the test of the ranking against measured
counts.

We decoded the 29{,}355 immune cells that the two
annotated studies label as lymphoid or myeloid, from control or eutopic
endometrium, from every split. Each cell was loaded with its own tissue,
platform, study and library size and encoded by the selected checkpoint.
Its latent was then moved one unit along the direction and decoded through
the same covariates, and the change in log expected count was averaged per
gene over the cells whose study measured it. A half-unit step gave the same
ranking.

The covariates must be the cell's own. The decoder multiplies each cell's
predicted profile by one number computed from the latent and the
covariates together. With every cell's real covariates that number was
negative, and with the covariates set to zero it changed sign for most
cells, which inverts every gene's relative level. Two checks therefore
preceded any ranking. First, every cell's encoded latent had to match the
stored latent to within 0.05, which confirms the checkpoint and the row
mapping. Second, six lymphoid and six myeloid markers had to decode the
right way round in labelled cells, at least five of six on each side, and
all twelve did.

The direction was the study-centred control-minus-patient difference of
donor means used throughout \ref{note:claims}, without removing its
component along the cell-family axes. Removing that component sent the
leading genes to transcripts that immune cells do not express. Genes that
immune cells express below 0.1 per 10{,}000 reads were excluded from the
ranking, because a decoded change on a gene the cells barely express has
no measured counterpart. Both choices were made after the ranking was seen
and before the measured-count test was run.

\textbf{Measured-count test.} The rule was fixed before the data existed.
For each donor with at least 50 immune cells, their raw counts were summed and scaled
as $\log(1 + 10^{4}\,x/n)$. In studies other than the two annotated ones,
immune cells were those with a measured PTPRC count above zero, a label
that never passes through the latent. For each of the fifteen genes
leading the ranking, the patient-minus-control difference was averaged
over studies, weighted by donors, and tested by permuting labels within
study 10{,}000 times. The claim required at least 11 of 15 genes higher in
patients and a sign-test $P < 0.05$, in held-out studies with at least
three patients and three controls. \texttt{li\_2025} was excluded because
its endometrium samples are follicular fluid (\ref{note:limitations}).
Overall, the decoding reads a direction through the model's own decoder
under the conditions it was trained on, and the test checks the result
against counts the model never produced.
